% Bibliographic data inherited verbatim from the current French authority.
\begin{thebibliography}{22}

\bibitem{I-1}
\oldpage[I]{214}
H.~\textsc{Cartan} et C.~\textsc{Chevalley},
\emph{Séminaire de l'École Normale Supérieure},
8\textsuperscript{e} année (1955-56), \emph{Géométrie algébrique}.

\bibitem{I-2}
H.~\textsc{Cartan} and S.~\textsc{Eilenberg},
\emph{Homological Algebra}, Princeton Math. Series (Princeton University
Press), 1956.

\bibitem{I-3}
W.~L.~\textsc{Chow} and J.~\textsc{Igusa},
\emph{Cohomology theory of varieties over rings},
\emph{Proc. Nat. Acad. Sci. U.S.A.}, t.~XLIV (1958), p.~1244-1248.

\bibitem{I-4}
R.~\textsc{Godement}, \emph{Théorie des faisceaux}, Actual. Scient. et Ind.,
n\textsuperscript{o}~1252, Paris (Hermann), 1958.

\bibitem{I-5}
H.~\textsc{Grauert},
\emph{Ein Theorem der analytischen Garbentheorie und die Modulräume komplexer
Strukturen}, \emph{Publ. Math. Inst. Hautes Études Scient.},
n\textsuperscript{o}~5, 1960.

\bibitem{I-6}
A.~\textsc{Grothendieck}, \emph{Sur quelques points d'algèbre homologique},
\emph{Tôhoku Math. Journ.}, t.~IX (1957), p.~119-221.

\bibitem{I-7}
A.~\textsc{Grothendieck},
\emph{Cohomology theory of abstract algebraic varieties},
\emph{Proc. Intern. Congress of Math.}, p.~103-118, Edinburgh (1958).

\bibitem{I-8}
A.~\textsc{Grothendieck},
\emph{Géométrie formelle et géométrie algébrique},
\emph{Séminaire Bourbaki}, 11\textsuperscript{e} année (1958-59), exposé 182.

\bibitem{I-9}
M.~\textsc{Nagata},
\emph{A general theory of algebraic geometry over Dedekind domains},
\emph{Amer. Math. Journ.}~: I, t.~LXXVIII, p.~78-116 (1956)~; II, t.~LXXX,
p.~382-420 (1958).

\bibitem{I-10}
D.~G.~\textsc{Northcott}, \emph{Ideal theory}, Cambridge Univ. Press, 1953.

\bibitem{I-11}
P.~\textsc{Samuel}, \emph{Commutative algebra} (Notes by D.~Herzig), Cornell
Univ., 1953.

\bibitem{I-12}
P.~\textsc{Samuel}, \emph{Algèbre locale}, Mém. Sci. Math.,
n\textsuperscript{o}~123, Paris, 1953.

\bibitem{I-13}
P.~\textsc{Samuel} and O.~\textsc{Zariski},
\emph{Commutative algebra}, 2 vol., New York (Van Nostrand), 1958-60.

\bibitem{I-14}
J.-P.~\textsc{Serre}, \emph{Faisceaux algébriques cohérents},
\emph{Ann. of Math.}, t.~LXI (1955), p.~197-278.

\bibitem{I-15}
J.-P.~\textsc{Serre},
\emph{Sur la cohomologie des variétés algébriques},
\emph{Journ. de Math.} (9), t.~XXXVI (1957), p.~1-16.

\bibitem{I-16}
J.-P.~\textsc{Serre},
\emph{Géométrie algébrique et géométrie analytique},
\emph{Ann. Inst. Fourier}, t.~VI (1955-56), p.~1-42.

\bibitem{I-17}
J.-P.~\textsc{Serre},
\emph{Sur la dimension homologique des anneaux et des modules noethériens},
\emph{Proc. Intern. Symp. on Alg. Number theory}, p.~176-189, Tokyo-Nikko,
1955.

\bibitem{I-18}
A.~\textsc{Weil}, \emph{Foundations of algebraic geometry}, Amer. Math. Soc.
Coll. Publ., n\textsuperscript{o}~29, 1946.

\bibitem{I-19}
A.~\textsc{Weil}, \emph{Numbers of solutions of equations in finite fields},
\emph{Bull. Amer. Math. Soc.}, t.~LV (1949), p.~497-508.

\bibitem{I-20}
O.~\textsc{Zariski},
\emph{Theory and applications of holomorphic functions on algebraic varieties
over arbitrary ground fields}, Mem. Amer. Math. Soc., n\textsuperscript{o}~5
(1951).

\bibitem{I-21}
O.~\textsc{Zariski}, \emph{A new proof of Hilbert's Nullstellensatz},
\emph{Bull. Amer. Math. Soc.}, t.~LIII (1947), p.~362-368.

\bibitem{I-22}
E.~\textsc{Kähler}, \emph{Geometria Arithmetica},
\emph{Ann. di Mat.} (4), t.~XLV (1958), p.~1-368.

\end{thebibliography}
